\section*{Основы теории механических волн}

\subsection*{1. Что такое волна — определение и смысл}

Волна — это пространственно-временное распространение возмущения в среде с переносом 
энергии и информации, но без макроскопического переноса вещества. Волны могут представлять 
собой колебания частиц среды (механические волны) или колебания физических полей 
(электромагнитные волны) \cite{openstax-waves, feynman}.

\subsection*{2. Математическая модель: одномерное волновое уравнение}

Для одномерного упругого стержня с малыми продольными смещениями \(u(x,t)\) получаем 
линейное волновое уравнение:
\[
\frac{\partial^2 u}{\partial t^2} = c^2 \frac{\partial^2 u}{\partial x^2},
\]
где \(c\) — скорость распространения волны. Для продольной волны:
\[
c = \sqrt{\frac{E}{\rho}},
\]
где \(E\) — модуль Юнга, \(\rho\) — плотность материала \cite{mit803}.

Общее решение в бесконечной однородной среде имеет вид:
\[
u(x,t) = f(x - ct) + g(x + ct),
\]
где \(f\) и \(g\) — произвольные функции профиля правой и левой волн \cite{mit803}.

\subsection*{3. Гармонические (плоские) волны}

Гармоническая волна задаётся выражением:
\[
u(x,t) = A \cos(kx - \omega t + \phi),
\]
где
\[
k = \frac{2\pi}{\lambda}, \quad \omega = 2\pi f.
\]
Дисперсионное соотношение \(\omega(k)\) связывает частоту и волновое число.  
Для недиспергирующей среды:
\[
\omega = ck,
\]
а фазовая и групповая скорости равны:
\[
v_{\text{ф}} = \frac{\omega}{k} = c, \quad
v_{\text{гр}} = \frac{d\omega}{dk} = c.
\]
Детальное обсуждение — в \cite{openstax-waves, mit803}.

\subsection*{4. Суперпозиция и стоячие волны}

Линейность волнового уравнения позволяет накладывать решения. Наложение двух 
противоположно направленных гармонических волн одинаковой частоты приводит к стоячим волнам:
\[
u(x,t) = 2A \sin(kx) \sin(\omega t).
\]
Подробный разбор — в \cite{openstax-waves, mit803}.

\subsection*{5. Дисперсия, фазовая и групповая скорость}

Дисперсия — это зависимость \(\omega(k)\). В диспергирующей среде разные гармоники волнового пакета движутся 
с разными скоростями, поэтому форма пакета изменяется со временем.

Фазовая скорость:
\[
v_p = \frac{\omega}{k},
\]
групповая скорость:
\[
v_g = \frac{d\omega}{dk}.
\]
Роль групповой скорости в передаче энергии подробно обсуждается в \cite{mit803}.

\subsection*{6. Энергия волны}

Механическая волна переносит кинетическую и потенциальную энергию. Для гармонической волны 
плотность энергии ~ \(A^2 \omega^2\) \cite{feynman, openstax-waves}.

\subsection*{7. Границы, отражение и преломление волн}

При переходе волны между двумя средами часть энергии отражается, часть проходит дальше. 
Коэффициенты отражения/передачи зависят от механических импедансов сред \cite{mit803, purdue-waves}.


\section*{Источники}

\begin{thebibliography}{9}

\bibitem{openstax-waves}
William Moebs, Samuel J. Ling, Jeff Sanny. \textit{University Physics, Volume 1 — Chapter 16 “Waves”}.  
OpenStax, 2016.  
Доступно онлайн: \url{https://assets.openstax.org/oscms-prodcms/media/documents/UniversityPhysicsVol1-WEB.pdf} \cite{openstax-waves}

\bibitem{feynman}
R. P. Feynman, R. B. Leighton, M. Sands. \textit{The Feynman Lectures on Physics, Vol. I}, Ch. 51 “Waves”.  
California Institute of Technology.  
Онлайн: \url{https://www.feynmanlectures.caltech.edu/I_51.html} \cite{feynman}

\bibitem{mit803}
MIT OpenCourseWare. \textit{8.03SC Physics III: Vibrations and Waves}.  
Лекции + конспекты. 2016.  
Ссылка на курс: \url{https://ocw.mit.edu/courses/8-03sc-physics-iii-vibrations-and-waves-fall-2016/} \cite{mit803}

\bibitem{purdue-waves}
Purdue University, Lecture Notes on Waves and Oscillations.  
Доступны на сайте Physics Department Purdue (примерно \url{https://www.physics.purdue.edu/}) \cite{purdue-waves}

\end{thebibliography}
